\document{article}

\title{LEGO Builder}

\section{Floating bricks}

Consider the following situation:

we have a model with two cluster of voxels that are far away each other (different heights).
For this situation, we would like to place bricks for both clusters as there's no way the two can be connected.

> Although alternative solutions were considered: for example, identifying such clusters in the model's grid and
> connect them with a straight line. Not feasible as it would lead to unstable structure.

Now, take the very first slice of the ''floating'' cluster.
Bricks placed for it should be valid placements.
Do we need a special treatment to accept those? 

Yes, if the placement isn't connecting any brick of the previous subslice, is invalid.

**But if the value of the previous subslice proximity map is below a threshold, then it's valid!**

\section{Reward function}

The program works by testing all the possible Placement that can be done for a 2d slice.
For every placement, a \textbf{reward function} is evaluated and the optimal placement is chosen.

The reward function evaluation is done leveraging GPU parallelism:
we can easily test all the possible, finite points, instead of performing an approximate search (for example, using Particle Swarm Optimization).

Given the following variables defined per placement $\rho$:

\begin{itemize}
    \item $a_{n} \in [0, 1]$ a number proportional to the number of sides with a neighbor over all placement sides
    \item $b_{n} \in ]0, 1]$ the ratio of the brick size over the maximum brick size
    \item $c_{n} \in [0, 1]$ a number proportional to the number of color spots covered by the placement
    \item $d_{n} \in [0, 1]$ a number proportional to the number of different bricks of the previous slice connected by the placement
\end{itemize}

The reward function is defined as:

\begin{equation}
    f(\rho) =
\end{equation}
